\documentclass[11pt]{article}
\usepackage[margin=1in]{geometry}
\usepackage{amsmath,amssymb}
\usepackage{array}
\usepackage{booktabs}
\usepackage{hyperref}
\usepackage{xcolor}
\usepackage{listings}

\lstset{
  basicstyle=\ttfamily\small,
  frame=single,
  breaklines=true,
  columns=fullflexible,
  keepspaces=true,
}

\newcommand{\phigold}{\varphi}
\newcommand{\Zphi}{\mathbb{Z}[\phigold]}

\title{SixMesh: Edges and Addresses in SixPhi Space}
\author{James A. Hinds (Jim) \and Claude (Opus 4.7)}
\date{2026-06-17}

\begin{document}
\maketitle

\begin{abstract}
A working chat log exploring the geometric structure of the
\texttt{space-struts} project's SixPhi coordinate system: how the six
icosahedral face-normal axes interact, the exact metric tensor \(G\), the
signed pattern of its entries, and the consequence that every edge in a
structure built from golden triangles carries a finite-alphabet
\emph{address} in SixPhi space.
\end{abstract}

\section{Setup: the project context}

\texttt{space-struts} is a live SvelteKit site exploring an exact
golden-ratio geometry system (``PhiBase'') and constructions of dodecahedral
and quasicrystal structures from golden triangles. Canonical modules:

\begin{itemize}
  \item \texttt{src/lib/coffee/phiBase.coffee} --- exact arithmetic in
        \(\Zphi\): \(P(p,n)/d = (p\phigold + n)/d\).
  \item \texttt{src/lib/coffee/sixPhiVector.coffee} --- the six-basis vectors.
  \item \texttt{src/lib/coffee/geoPhi.coffee} --- the geometry engine
        (\texttt{GeoPhi} class, the 15 mirror planes, neighbor star,
        golden-apex completion).
\end{itemize}

\section{The six basis vectors}

The six SixPhi basis directions are the face-normals of the icosahedron
(the 5-fold axes \(A\)--\(F\)):

\[
\begin{aligned}
b_0 &= (\phigold,\,0,\,1)   & b_1 &= (\phigold,\,0,\,-1) \\
b_2 &= (0,\,1,\,\phigold)   & b_3 &= (0,\,-1,\,\phigold) \\
b_4 &= (1,\,\phigold,\,0)   & b_5 &= (-1,\,\phigold,\,0)
\end{aligned}
\]

with antipodal pairing \((0\!\leftrightarrow\!1),\ (2\!\leftrightarrow\!3),\
(4\!\leftrightarrow\!5)\).

\section{The observation: every step has a fixed shadow}

\textbf{Jim's observation.}\quad
\emph{Each movement of \(1\) or \(\phigold\) along an axis adds some fixed
result along that axis, and the other five SixPhi axes would each get some
fixed result for that movement.}

This is exactly the Gram matrix
\(G_{ij} = b_i \cdot b_j\). Computing exactly:
\begin{itemize}
  \item Diagonal: \(b_k \cdot b_k = \phigold^2 + 1 = \phigold + 2 = P(1,2)\).
        (The same \(\phigold+2\) whose algebraic norm is \(5\), giving the
        reflection denominator.)
  \item Off-diagonal: \(b_i \cdot b_j = \pm\phigold = P(\pm 1, 0)\) for every
        \(i \neq j\).
\end{itemize}

So a unit step along axis \(k\) contributes \(\pm\phigold\) to each other
axis; a step of \(\phigold\) along \(k\) contributes
\(\pm\phigold^2 = \pm(\phigold+1)\); and the move registers \(\phigold+2\)
on \(k\) itself.

\section{The full signed \texorpdfstring{$G$}{G} table}

\[
G \;=\;
\begin{array}{c|cccccc}
        & b_0      & b_1      & b_2      & b_3      & b_4      & b_5      \\ \hline
b_0     & \phigold{+}2 & +\phigold    & +\phigold    & +\phigold    & +\phigold    & -\phigold    \\
b_1     & +\phigold    & \phigold{+}2 & -\phigold    & -\phigold    & +\phigold    & -\phigold    \\
b_2     & +\phigold    & -\phigold    & \phigold{+}2 & +\phigold    & +\phigold    & +\phigold    \\
b_3     & +\phigold    & -\phigold    & +\phigold    & \phigold{+}2 & -\phigold    & -\phigold    \\
b_4     & +\phigold    & +\phigold    & +\phigold    & -\phigold    & \phigold{+}2 & +\phigold    \\
b_5     & -\phigold    & -\phigold    & +\phigold    & -\phigold    & +\phigold    & \phigold{+}2 \\
\end{array}
\]

Stripped to its off-diagonal sign pattern:

\[
\begin{array}{c|cccccc}
   & 0 & 1 & 2 & 3 & 4 & 5 \\ \hline
0  & \cdot & + & + & + & + & - \\
1  & +     & \cdot & - & - & + & - \\
2  & +     & -     & \cdot & + & + & + \\
3  & +     & -     & +     & \cdot & - & - \\
4  & +     & +     & +     & -     & \cdot & + \\
5  & -     & -     & +     & -     & +     & \cdot \\
\end{array}
\]

\paragraph{Structural notes.}
\begin{itemize}
  \item Every row has exactly \textbf{4 plus} and \textbf{1 minus} off the
        diagonal.
  \item Per-row off-diagonal sums alternate by antipodal pair:
        \(r_0,r_2,r_4 = +3\phigold\); \(r_1,r_3,r_5 = -\phigold\).
  \item The lone minus in row \(k\) traces a perfect matching across the
        three antipodal pairs: \(0\!\leftrightarrow\!5,\
        1\!\leftrightarrow\!2,\ 3\!\leftrightarrow\!4\) (three
        transpositions).
  \item Because every off-diagonal is \(\pm\phigold\), the matrix has only
        two distinct entries up to sign --- this is why pure six-basis
        translation stays in \(\Zphi\) and never leaks into
        \(\tfrac{1}{5}\Zphi\).
\end{itemize}

\section{Edges as SixPhi addresses}

\textbf{Jim's next move.}\quad
\emph{Imagine a structure made of our triangles, always laid out in SixPhi
space, each edge connecting two vertices, each with their coordinates in
SixPhi. Each edge would then have an ``address'' in SixPhi space.}

Define the edge address as the displacement
\(\Delta = v_2 - v_1\) in SixPhi coordinates. Two edges are translation
equivalent iff they share an address (up to sign). Every edge factors into:

\begin{enumerate}
  \item \textbf{Direction class} --- one of 30 unsigned neighbor-star
        directions.
  \item \textbf{Length class} --- short \(s = 2/\phigold\) or long
        \(L = 2\) (the only two edge lengths a golden triangle or gnomon
        admits).
  \item \textbf{Sign} --- orientation along the direction.
\end{enumerate}

This alphabet is already in the codebase: \texttt{GeoPhi.neighborStar} is
the \emph{finite alphabet of edge addresses}, \(60 = 30 + 30\): the two
\(I_h\) orbits of the golden edge vectors (short \(s\) and long \(L\)).

\subsection{Consequences}

\begin{description}
  \item[Compression.] A structure reduces to a vertex list plus an edge
        list of triples \((\text{idx}_a, \text{idx}_b,
        \text{neighborStarIndex})\) --- displacement is redundant with the
        address.
  \item[Cheap invariants.] ``How many short edges along direction \(k\)?''
        is a histogram over a 60-bin alphabet. Two structures are
        translation-congruent iff edge-address histograms match (necessary;
        sufficient with one vertex fixed).
  \item[Closure.] A triangle \((a,b,c)\) is golden iff its three edge
        addresses sum to zero in SixPhi \emph{and} length classes form
        \(\{s,L,L\}\) (triangle) or \(\{s,s,L\}\) (gnomon).
        \texttt{goldenApexCandidates} essentially solves: given two
        addresses, what third address closes a legal golden tile?
  \item[Plane signatures lift to edges.] Each neighbor-star direction lies
        in a subset of the 15 mirror planes (via the exact
        \texttt{planesContaining} test), so every edge inherits a plane
        signature; a structure inherits the intersection of its edges'
        signatures --- a fingerprint of which \(I_h\) symmetries it
        respects.
  \item[\(\Zphi\) vs.\ \(\tfrac{1}{5}\Zphi\).] Pure edge-star translation
        keeps addresses in \(\Zphi\). Reflection across any of the 15
        mirror planes drops the new edges into \(\tfrac{1}{5}\Zphi\) (norm
        of \(\phigold+2\) is 5). The address itself tells you which regime
        you're in.
\end{description}

\section{Open directions}

\begin{itemize}
  \item Enumerate the 60 neighbor-star addresses as a labeled table
        (direction \(\times\) length \(\times\) mirror-plane signature).
  \item Sketch the edge-list data structure as a small CoffeeScript helper
        in \texttt{geoPhi.coffee}.
  \item Work out the triangle-closure algebra: given addresses
        \((\alpha,\beta)\), enumerate the closing \(\gamma\) with
        \(\alpha+\beta+\gamma=0\) and lengths in \(\{s,L,L\}\) or
        \(\{s,s,L\}\).
\end{itemize}

\end{document}
